\documentclass{article}
\usepackage{amsmath, amssymb, graphicx}

\begin{document}

\title{Modelagem da Probabilidade de Ausência Baseada na Idade da Mãe}
\author{Mirna}
\date{}
\maketitle

\section{Fundamentação Estatística: Probabilidade de Ausência}

O objetivo é modelar a probabilidade condicional de ausência da variável \texttt{IDADEPAI} com base na \texttt{IDADEMAE}. Isso significa que, à medida que a idade da mãe diminui, a probabilidade de ausência do dado aumenta.

Formalmente, queremos modelar a seguinte probabilidade condicional:

\[
P(\text{IDADEPAI ausente} \mid \text{IDADEMAE})
\]

\subsection{Função Logística}

Uma forma comum de modelar essa relação é por meio da função logística:

\[
P(\text{IDADEPAI ausente}) = \frac{1}{1 + e^{-\beta (\text{IDADEMAE} - \mu)}}
\]

onde:

\begin{itemize}
    \item \( \beta \) controla a taxa de variação da probabilidade.
    \item \( \mu \) representa a idade média das mães na população.
    \item \( e \) é a base do logaritmo natural (aproximadamente 2.718).
\end{itemize}

A função logística é útil porque:

\begin{enumerate}
    \item Gera probabilidades no intervalo [0,1].
    \item Modela efeitos suaves: a transição da ausência para a presença ocorre de maneira gradual.
    \item É amplamente utilizada em modelos de regressão logística.
\end{enumerate}

No código em R, isso pode ser representado por:

\begin{verbatim}
dt_simulado[, prob_missing := 1 / (1 + exp(0.3 * (IDADEMAE - mean(IDADEMAE, na.rm = TRUE))))]
\end{verbatim}

\section{Interpretação do Modelo}

Se representarmos graficamente essa função logística, podemos observar que:

\begin{itemize}
    \item Mães mais jovens têm maior probabilidade de ausência de \texttt{IDADEPAI}.
    \item Se \texttt{IDADEMAE} for muito abaixo da média, \( P(\text{missing}) \) tende a 1.
    \item Se \texttt{IDADEMAE} for acima da média, \( P(\text{missing}) \) tende a 0.
\end{itemize}

A tabela a seguir mostra a variação da probabilidade de ausência para diferentes valores de \(\beta\):

\begin{table}[h]
    \centering
    \begin{tabular}{|c|c|c|c|}
        \hline
        \textbf{IDADEMAE}  & \(\beta = 0.1\) (suave) & \(\beta = 0.3\) (médio) & \(\beta = 0.5\) (íngreme) \\
        \hline
        18 anos  & 80\% missing  & 90\% missing  & 95\% missing  \\
        25 anos  & 50\% missing  & 60\% missing  & 75\% missing  \\
        35 anos  & 10\% missing  & 20\% missing  & 40\% missing  \\
        \hline
    \end{tabular}
    \caption{Variação da probabilidade de ausência conforme \texttt{IDADEMAE} e \(\beta\)}
\end{table}

A escolha do coeficiente \(\beta\) afeta a rapidez com que a transição ocorre:

\begin{itemize}
    \item Se \(\beta\) for pequeno (\(0.1\)), a mudança de probabilidade é suave.
    \item Se \(\beta\) for grande (\(0.5\)), a mudança é mais abrupta.
\end{itemize}

\section{Relação com Modelos Estatísticos}

Esse modelo se baseia em duas abordagens estatísticas:

\subsection{Regressão Logística}

A regressão logística é um modelo amplamente usado para modelar variáveis binárias. A forma geral do modelo logit é:

\[
\log \left( \frac{P(Y=1)}{1 - P(Y=1)} \right) = \beta_0 + \beta_1 X
\]

onde \( Y \) é a variável indicadora da ausência de \texttt{IDADEPAI} e \( X = \text{IDADEMAE} \).

\subsection{Modelos de Risco Proporcional}

Outra abordagem possível é o modelo de risco proporcional de Cox, amplamente usado em análise de sobrevivência. A função logística pode ser interpretada como um modelo de risco para probabilidade de eventos.

\section{Ajustando os Parâmetros do Modelo}

Podemos ajustar \(\beta\) empiricamente usando um modelo de regressão logística em R:

\begin{verbatim}
glm(is.na(IDADEPAI) ~ IDADEMAE, family = binomial, data = df)
\end{verbatim}

Isso permite estimar automaticamente o coeficiente \(\beta\) com base nos dados reais.

\section{Conclusão}

\begin{itemize}
    \item Utilizamos uma função logística para modelar a ausência de \texttt{IDADEPAI} como uma função de \texttt{IDADEMAE}.
    \item A relação capturada reflete um mecanismo do tipo MAR (Missing at Random).
    \item O coeficiente \(\beta\) pode ser ajustado para controlar a intensidade do efeito.
\end{itemize}




\section{Modelagem de Ausência de Dados e Avaliação de Métodos de Imputação}
\subsection{Introdução}

O objetivo deste estudo é avaliar diferentes métodos de imputação para dados ausentes na variável \texttt{IDADEPAI}. Para isso, utilizamos simulações com três mecanismos de ausência de dados: MCAR (Missing Completely at Random), MAR (Missing at Random) e MNAR (Missing Not at Random). Os dados imputados são avaliados por meio de métricas estatísticas, comparando-os com os valores verdadeiros da população completa.

\section{Cenários de Ausência de Dados}

Para cada simulação, removemos uma proporção \( p \) dos dados de acordo com um dos seguintes mecanismos:

\begin{itemize}
    \item \textbf{MCAR:} A ausência ocorre completamente ao acaso, sem relação com outras variáveis.
    \item \textbf{MAR:} A ausência está relacionada à idade da mãe (\texttt{IDADEMAE}), onde menores valores de \texttt{IDADEMAE} aumentam a probabilidade de \texttt{IDADEPAI} estar ausente. Esse mecanismo pode ser modelado como uma regressão logística:

    \begin{equation}
        P(\texttt{IDADEPAI} = \text{NA} \mid \texttt{IDADEMAE}) = \frac{1}{1 + e^{\beta_0 + \beta_1 \texttt{IDADEMAE}}}
    \end{equation}

    \item \textbf{MNAR:} A ausência depende diretamente do valor da própria variável \texttt{IDADEPAI}, sendo mais frequente para pais mais jovens:

    \begin{equation}
        P(\texttt{IDADEPAI} = \text{NA} \mid \texttt{IDADEPAI} < 30) > P(\texttt{IDADEPAI} = \text{NA} \mid \texttt{IDADEPAI} \geq 30)
    \end{equation}
\end{itemize}

\section{Métodos de Imputação}

Os métodos utilizados para preencher os valores ausentes incluem:

\begin{itemize}
    \item \textbf{Casos completos:} Apenas os registros sem valores ausentes são utilizados na análise.
    \item \textbf{Imputação por Predictive Mean Matching (PMM):} Método baseado em vizinhos mais próximos que preserva a distribuição original.
    \item \textbf{Imputação por Árvores de Decisão (CART):} Utiliza árvores de regressão para prever valores ausentes.
    \item \textbf{Imputação por Random Forest:} Utiliza florestas aleatórias para prever valores ausentes.
\end{itemize}

\section{Métricas de Avaliação}

A qualidade das imputações é avaliada por meio das seguintes métricas:

\subsection{Erro Quadrático Médio (RMSE)}
Mede a diferença média entre os valores verdadeiros (\( y_i \)) e os imputados (\( \hat{y}_i \)):

\begin{equation}
    RMSE = \sqrt{\frac{1}{n} \sum_{i=1}^{n} (y_i - \hat{y}_i)^2}
\end{equation}

\subsection{Viés Relativo (RB)}
Compara a média dos valores imputados com a média dos valores verdadeiros:

\begin{equation}
    RB = \frac{\bar{\hat{y}} - \bar{y}}{\bar{y}}
\end{equation}

\subsection{Proporção de Viés (PB)}
Relaciona a diferença entre as médias imputadas e verdadeiras ao desvio padrão dos valores verdadeiros:

\begin{equation}
    PB = \frac{\bar{\hat{y}} - \bar{y}}{s_y}
\end{equation}

\subsection{Erro Absoluto Médio (MAE)}
Mede a magnitude média dos erros absolutos entre os valores imputados e os valores verdadeiros:

\begin{equation}
    MAE = \frac{1}{n} \sum_{i=1}^{n} |y_i - \hat{y}_i|
\end{equation}

\subsection{Erro Percentual Absoluto Médio (MAPE)}
Expressa o erro absoluto médio como uma porcentagem dos valores verdadeiros:

\begin{equation}
    MAPE = \frac{1}{n} \sum_{i=1}^{n} \left| \frac{y_i - \hat{y}_i}{y_i} \right| \times 100
\end{equation}

\section{Conclusão}

Através das simulações, analisamos o impacto de diferentes mecanismos de ausência de dados e métodos de imputação. A escolha do método de imputação ideal depende do mecanismo subjacente de ausência e do objetivo da análise. Em geral, métodos baseados em machine learning, como Random Forest, tendem a apresentar menor viés e erro quadrático médio.

\end{document}
